\section{付録}
\subsection{高次項を含むハミルトニアンの厳密な展開とRWAの詳細}
2.1.1節で導入したトランスモンハミルトニアンの近似について、第4次までのテイラー展開の詳細な過程を示す。
演算子 $\hat{\phi}$ を生成消滅演算子で表現すると、

\begin{equation}
  \hat{\phi} = \left(\frac{8E_C}{E_J}\right)^{1/4} \frac{\hat{a} + \hat{a}^\dagger}{\sqrt{2}}
\end{equation}

これをポテンシャル項 $-E_J \cos(\hat{\phi})$ に代入して展開する。

\begin{equation}
  -E_J \left( 1 - \frac{\hat{\phi}^2}{2!} + \frac{\hat{\phi}^4}{4!} - \dots \right)
\end{equation}

$\hat{\phi}^4$ の項は、正規順序化（Normal Ordering）を行うことで、非調和性をもたらす項 $\hat{a}^\dagger \hat{a}^\dagger \hat{a} \hat{a}$ と、周波数シフトに寄与する項に分離される。ここで回転波近似（RWA）を適用し、粒子数を保存しない項（例：$\hat{a}^4$ や $\hat{a}^{\dagger 4}$）を削除することで、有効ハミルトニアンが導出される。

\subsection{断熱定理に基づくCZ条件の数学的記述}
断熱的な軌道 $|11\rangle \rightarrow \overline{|11\rangle} \rightarrow |11\rangle$ において、非断熱誤差を最小化するための断熱条件は以下の不等式で与えられる。

\begin{equation}
  \left| \frac{\langle 02 | \partial_t \hat{H} | 11 \rangle}{(E_{11} - E_{02})^2} \right| \ll 1
\end{equation}

この条件を満たしつつ、ゲート操作時間 $\tau$ をデコヒーレンス時間（$T_1, T_2$）よりも十分に短くするために、微分の連続性を担保したエンベロープ形状が要求される。
